%! Author = lili
%! Date = 7/15/26

\newglossaryentry{testkarte}{
    name={test},
    description={todo},
    type=dinge
}

% =====================================================================
%  Karteikarten  Vorlesung 4
% =====================================================================

\newglossaryentry{dadp}{
    name={Was sind beispiele für DNA-abhängige DNA-Polymerasen und wo werden sie eingesetzt?},
    description={Replikation, Taq-Polymerase},
    type=dinge
}

\newglossaryentry{darp}{
    name={Was sind beispiele für DNA-abhängige RNA-Polymerasen und wo werden sie eingesetzt?},
    description={Transkription, Primase bei Replikation},
    type=dinge
}

\newglossaryentry{radp}{
    name={Wo werden RNA-abhängige DNA-Polymerasen eingesetzt?},
    description={reverse Transkription (RNA-Viren)},
    type=dinge
}

\newglossaryentry{rarp}{
    name={Wo werden RNA-abhängige RNA-Polymerasen eingesetzt?},
    description={post-transcriptional gene silencing
    (PTGS), "silencing"},
    type=dinge
}


\newglossaryentry{transkrtranslkopplung}{
    name={Worin unterscheidet sich die zeitliche Kopplung von Transkription und Translation bei Pro- und Eukaryoten?},
    description={Bei Prokaryoten läuft die Translation (fast) gleichzeitig mit der Transkription ab (gekoppelt, am selben Ort). Bei Eukaryoten sind beide Prozesse räumlich und zeitlich getrennt: Transkription im Zellkern, Translation im Cytoplasma.},
    type=dinge
}

\newglossaryentry{rnapolkeinprimer}{
    name={Warum brauchen RNA-Polymerasen im Gegensatz zu DNA-Polymerasen keinen Primer?},
    description={RNA-Polymerasen können die Synthese de novo starten -- sie initiieren die RNA-Kette ohne vorgelegtes freies 3'-OH-Ende. DNA-Polymerasen dagegen können nur an ein bestehendes 3'-OH-Ende (Primer) anknüpfen.},
    type=dinge
}

\newglossaryentry{welcherstrangvorlage}{
    name={Welcher Strang der DNA-Doppelhelix dient als Vorlage für die Transkription?},
    description={Je nach Gen kann mal der eine, mal der andere Strang als Matrize (Vorlage) verwendet werden -- es ist nicht immer derselbe Strang. Katalysiert wird die Synthese von einer DNA-abhängigen RNA-Polymerase.},
    type=dinge
}

\newglossaryentry{transkriptionsgleichung}{
    name={Wie lautet die (vereinfachte) Reaktionsgleichung der Transkription?},
    description={DNA + n NTP  --(RNA-Polymerase)-->  DNA + RNA + n PP. Die DNA bleibt erhalten (dient nur als Matrize), pro eingebautem Nukleotid wird ein Pyrophosphat (PP) freigesetzt.},
    type=dinge
}
\newglossaryentry{codingvsncrna}{
    name={Wie unterscheidet man RNA-Klassen grob, und wie verändert sich der Anteil nicht-kodierender RNA (ncRNA) mit der Komplexität?},
    description={Man unterscheidet protein-kodierende (coding) RNA und nicht-kodierende RNA (ncRNA). Der Anteil der ncRNA steigt mit zunehmender Komplexität eines Organismus. Es werden laufend neue RNA-Klassen entdeckt (benannt nach Funktion oder Lokalisation).},
    type=dinge
}

\newglossaryentry{polklassenzuordnung}{
    name={Welche RNA-Klassen synthetisieren RNA-Pol I, II und III jeweils?},
    description={Pol I: die großen rRNAs (28S, 18S, 5.8S). Pol II: hnRNA/Primärtranskripte, mRNA, viele miRNAs, einige snRNAs. Pol III: tRNA, 5S-rRNA, einige snRNAs (z.B. U6).},
    type=dinge
}

\newglossaryentry{kleinernazuordnung}{
    name={Von welcher Polymerase werden snRNA, miRNA und siRNA transkribiert?},
    description={snRNA: teils Pol II, teils Pol III (z.B. U6 durch Pol III). miRNA (20--22 nt): Pol II. siRNA (24 nt, v.a. Pflanzen): Pol IV und Pol V, unter Beteiligung einer RdRP.},
    type=dinge
}
\newglossaryentry{rnapolaufbau}{
    name={Wie ist eine eukaryotische RNA-Polymerase grob aufgebaut?},
    description={Als großer Multiprotein-Komplex aus mehreren Untereinheiten (bis ca. 12) mit konserviertem katalytischem Kern. Besonderheit der Pol II: eine C-terminale Domäne (CTD) an der größten Untereinheit, die während der Initiation/Elongation phosphoryliert wird.},
    type=dinge
}

\newglossaryentry{nukleolusaufbau}{
    name={Wo werden die eukaryotischen Ribosomen zusammengebaut, und wie entsteht der Nukleolus?},
    description={Die Ribosomen werden im Nukleolus zusammengebaut. Der Nukleolus geht vom Nukleolus-Organisator aus, der aus der NOR (nucleolus organizer region) entsteht.},
    type=dinge
}

\newglossaryentry{rdnasichtbar}{
    name={Was ist mit der "sichtbaren Genexpression" der rDNA gemeint?},
    description={Bei der rDNA kann man die Genexpression direkt "sehen" (z.B. als "Weihnachtsbaum"-Struktur im Elektronenmikroskop), weil das Genprodukt selbst RNA (rRNA) ist und nicht erst in ein Protein übersetzt werden muss.},
    type=dinge
}

\newglossaryentry{mrnaenden}{
    name={Was muss man über die Enden der eukaryotischen mRNA wissen (Vergleich zu prokaryotischer)?},
    description={Beide Enden sind modifiziert: am 5'-Ende eine Cap-Struktur (m7Gppp), am 3'-Ende ein Poly(A)-Schwanz. Beide schützen die mRNA vor Abbau und fördern die Translation. Prokaryotische mRNA besitzt diese Modifikationen nicht.},
    type=dinge
}

\newglossaryentry{mrnaelemente}{
    name={Nenne die funktionellen Elemente der eukaryotischen mRNA (5'->3') und ihre Funktion},
    description={5'-Cap (m7G/GTP): Schutz vor Abbau, Bindung des Ribosoms. Leader (5'-UTR): Ribosomenbindung. AUG: Startcodon (Translationsbeginn mit Methionin). ORF: kodiert die Aminosäure-Abfolge. Stopcodon (UGA/UAA/UAG): keine beladene tRNA vorhanden. Trailer (3'-UTR) mit AAUAAA: Signal für Schnitt und Polyadenylierung. Poly(A)-Schwanz: Stabilisierung der mRNA.},
    type=dinge
}

\newglossaryentry{polyawissen}{
    name={Was muss man über den Poly(A)-Schwanz wissen (Länge, Herkunft, Enzym, Bindeprotein, Funktion)?},
    description={Etwa 200 nt lang und NICHT im Genom kodiert. Wird nach dem Abschneiden des Transkripts von der Poly(A)-Polymerase angehängt und von PABP (poly(A)-binding protein) gebunden. Er stabilisiert die mRNA und erleichtert (mit PABP) die Translation. Bei stabilen, effizient translatierten mRNAs kann er auch kurz sein (~30 Adenosinreste).},
    type=dinge
}

\newglossaryentry{deadenylierungfolge}{
    name={Was ist Deadenylierung und welche Folge hat sie?},
    description={Die Entfernung bzw. Verkürzung des Poly(A)-Schwanzes. Dadurch löst sich PABP ab, was den Abbau der mRNA einleitet.},
    type=dinge
}

\newglossaryentry{mrnahalbwertszeit}{
    name={Wodurch wird die Halbwertszeit einer mRNA reguliert (trotz Poly(A))?},
    description={Durch spezielle Signale in der mRNA-Sequenz, oft im Trailer (3'-UTR). Eine mRNA kann trotz vorhandenem Poly(A)-Schwanz sehr kurzlebig sein.},
    type=dinge
}

\newglossaryentry{promotorallgemein}{
    name={Was bestimmt der Promotor, und wie unterscheidet er sich bei Eu- vs. Prokaryoten?},
    description={Der Promotor bestimmt den Start (und die Richtung) der Transkription; er ist durch spezifische DNA-Sequenzen ("Boxen") charakterisiert. Bei Eukaryoten ist er komplexer organisiert als bei Prokaryoten. Jede RNA-Polymerase erkennt ihre eigenen, spezifischen Promotorsequenzen, die ihre Zielgene definieren.},
    type=dinge
}

\newglossaryentry{promotorpol1}{
    name={Wie ist der Promotor der RNA-Pol I aufgebaut und was transkribiert sie?},
    description={Pol I transkribiert die 28S-, 18S- und 5.8S-rRNA. Ihr Promotor liegt extern (upstream, vor dem Transkriptionsstart) und besteht aus dem Core-Promotor + UCE (upstream control element).},
    type=dinge
}

\newglossaryentry{promotorpol2}{
    name={Wie ist der Promotor der RNA-Pol II aufgebaut und was transkribiert sie?},
    description={Pol II transkribiert mRNA, hnRNA, viele miRNAs und einige snRNAs. Ihr Promotor liegt extern (upstream): häufig eine TATA-Box (~ -25) und/oder ein Inr, oft zusätzlich CAAT-Box, GC-Box und weitere regulatorische Elemente.},
    type=dinge
}

\newglossaryentry{promotorpol3}{
    name={Wie ist der Promotor der RNA-Pol III aufgebaut und was transkribiert sie?},
    description={Pol III transkribiert tRNA, 5S-rRNA und einige snRNAs. Ihr Promotor liegt meist intern -- häufig innerhalb des transkribierten Gens (interner Promotor), z.B. Box A und Box B.},
    type=dinge
}

\newglossaryentry{promotorenvergleich}{
    name={Wie unterscheiden sich die Promotoren der drei RNA-Polymerasen grundsätzlich?},
    description={Pol I und Pol II nutzen externe (upstream) Promotoren, die vor dem Gen liegen. Pol III nutzt meist einen internen Promotor innerhalb des transkribierten Gens. Zudem hat jede Polymerase eigene charakteristische Promotorelemente (Pol I: Core + UCE; Pol II: TATA/Inr etc.; Pol III: Box A/B).},
    type=dinge
}

\newglossaryentry{tbpbesonderheit}{
    name={Was ist am TBP (TATA-Box-bindendes Protein) besonders (Bindung, Struktur, Rolle)?},
    description={TBP bindet in der kleinen Grube (minor groove) der DNA und besteht aus zwei Domänen mit ca. 40 \% Identität. Es ist ein zentraler Faktor der Promotorerkennung bei ALLEN drei RNA-Polymerasen.},
    type=dinge
}

\newglossaryentry{promotorelementekategorien}{
    name={Ordne die regulatorischen DNA-Elemente eukaryotischer Promotoren in drei Kategorien ein (mit Beispielen)},
    description={Core-Promotor: TATA-Box, Inr, BRE, DPE, MTE. Proximale Promotorelemente: CAAT-Box, GC-Box. Distale regulatorische Elemente: UAS, URS, RE, Enhancer, Silencer.},
    type=dinge
}

\newglossaryentry{basaletranskrivitroinvivo}{
    name={Was reicht für die basale Transkription aus, und warum ist sie in vivo unterdrückt?},
    description={In vielen experimentellen Systemen reicht ein Promotor mit nur Inr und TATA-Box für die basale Transkription. In vitro läuft sie ab; in vivo ist sie jedoch stark moduliert/unterdrückt, weil ohne regulatorische Faktoren nur eine Grundaktivität entsteht, die in der Zelle durch Chromatin und Regulatoren kontrolliert wird.},
    type=dinge
}

\newglossaryentry{corepromotorelementepol2}{
    name={Welche Core-Promotor-Elemente der RNA-Pol II kommen in unterschiedlichen Kombinationen vor?},
    description={BRE, TATA-Box, Inr und DPE. Sie treten in verschiedenen Kombinationen auf und bestimmen gemeinsam die Transkriptionskontrolle der Pol II.},
    type=dinge
}

\newglossaryentry{distale2d3d}{
    name={Wie können weit entfernte regulatorische Elemente (z.B. Enhancer) auf ein Gen wirken (2D vs. 3D)?},
    description={In der linearen 2D-Darstellung sind sie weit vom Gen entfernt. In der 3D-Organisation des Chromatins werden sie durch Schleifenbildung (looping) räumlich an den Promotor herangebracht und können so direkt auf die Transkription wirken.},
    type=dinge
}

\newglossaryentry{gtfvsrtf}{
    name={Was ist der Unterschied zwischen allgemeinen (GTFs) und regulatorischen Transkriptionsfaktoren (R-TFs)?},
    description={Allgemeine TFs (GTFs) sind für die basale Transkription notwendig und bilden mit der Pol II den Initiationskomplex. Regulatorische TFs (R-TFs) beeinflussen die Transkriptionsrate spezifisch (aktivierend oder reprimierend).},
    type=dinge
}

\newglossaryentry{allgemeinetfpol2}{
    name={Nenne die allgemeinen Transkriptionsfaktoren der RNA-Pol II und ihre Hauptfunktion},
    description={TFIID: mit TBP Erkennung der TATA-Box, mit TAFs Erkennung TATA-freier Promotoren und regulatorische Funktionen. TFIIA: stabilisiert die TFIID-Bindung. TFIIB: bindet die Pol II. TFIIF: führt die Pol II an den Promotor heran. TFIIE: führt TFIIH heran und reguliert dessen Aktivität. TFIIH: Helikase (Entwindung des Promotors) und Protein-Kinase (Phosphorylierung der CTD der Pol II).},
    type=dinge
}

\newglossaryentry{tfiihfunktion}{
    name={Welche enzymatischen Aktivitäten hat TFIIH und welche Untereinheiten gehören dazu?},
    description={Zwei DNA-Helikasen zur Promotor-Entwindung: XPB (3'->5') und XPD (5'->3'). Außerdem eine cyclinabhängige Protein-Kinase CDK7 (mit Cyclin H als regulatorischer Untereinheit), die die CTD der Pol II phosphoryliert.},
    type=dinge
}

\newglossaryentry{picII}{
    name={Was ist der Präinitiationskomplex (PIC) und woraus besteht er?},
    description={Der Protein-Komplex, der sich am Core-Promotor der RNA-Pol II bildet, bevor die Transkription startet. Er besteht aus der RNA-Pol II und den allgemeinen Transkriptionsfaktoren TFIIA, TFIIB, TFIID, TFIIE, TFIIF und TFIIH.},
    type=dinge
}


\newglossaryentry{elongationsschritte}{
    name={In welche drei Schritte lässt sich die Elongationsphase unterteilen, und wie viele Faktoren sind beteiligt?},
    description={i) Erkennung des Codons durch die Aminoacyl-tRNA; ii) Ausbildung der Peptidbindung zwischen zwei Aminosäuren; iii) Verschiebung (Translokation) des tRNA-mRNA-Komplexes. An diesen Prozessen sind vier eEFs (eukaryotische Elongationsfaktoren) beteiligt.},
    type=dinge
}

\newglossaryentry{translationstermination}{
    name={Wie läuft die Termination (der Translation) über Stopcodons und Release-Faktoren ab?},
    description={Gelangt ein Stopcodon (UAA, UAG oder UGA) in die A-Position des Ribosoms, kann keine tRNA binden (es gibt keine tRNA, die diese Tripletts erkennt). Die Release-Faktoren RF1 und RF2 erkennen die Stopcodons; ihre Bindung an die A-Stelle führt zur Abspaltung der Polypeptidkette von der tRNA in der P-Position, und das Ribosom zerfällt in seine Untereinheiten.},
    type=dinge
}

\newglossaryentry{transkriptionstermination}{
    name={Wie unterscheidet sich die Termination der Transkription bei RNA-Pol I, II und III?},
    description={Pol I: klares Terminationssignal (18-bp-Sequenz + Hilfsfaktor). Pol II: KEIN klares Signal -> "Torpedo"-Modell. Pol III: klares Signal ((G/C)n UUUU (G/C)n). Die Einzelheiten sind bei Eukaryoten wenig bekannt; alle drei Polymerasen nutzen unterschiedliche Mechanismen.},
    type=dinge
}

\newglossaryentry{mrnareifung}{
    name={Welche drei Schritte umfasst die Reifung des Primärtranskripts zur fertigen mRNA in Eukaryoten?},
    description={1) Anfügen der m7Gppp-Kappe (Cap) am 5'-Ende. 2) Spaltung und Polyadenylierung des 3'-Endes. 3) Spleißen (Entfernen der Introns).},
    type=dinge
}

\newglossaryentry{histonmrna}{
    name={Warum ist Histon-mRNA eine Ausnahme, und wie wird ihr 3'-Ende gebildet?},
    description={Canonical Histon-mRNA hat KEINEN Poly(A)-Schwanz; ihre Expression ist an die Replikation gekoppelt und zellzyklusreguliert. Das 3'-Ende erfordert einen konservierten Hairpin (stem-loop) und das HDE (histone downstream element), das mit der U7-snRNA hybridisiert. Die Schnittreaktion braucht SLBP (stem-loop binding protein) und U7-snRNA; den Schnitt führt CPSF aus (das auch bei Poly(A)-mRNAs schneidet).},
    type=dinge
}

\newglossaryentry{regulationbakterieneuk}{
    name={Vergleiche Transkription/Genregulation in Bakterien und Eukaryoten (wesentliche Unterschiede)},
    description={Bakterien: kein Zellkern; polycistronische Transkripte; Operons; eine RNA-Pol mit alternativen Sigma-Faktoren; Consensus-Promotoren; Shine-Dalgarno-Sequenz; transkriptionelle Attenuation. Eukaryoten: Zellkern; monocistronische Transkripte; keine Operons; mehrere RNA-Polymerasen für verschiedene Genklassen; sehr diverse Promotoren; Cap statt Shine-Dalgarno; nukleäres Spleißen; Enhancer, Chromatinstruktur sowie generelle und spezielle TFs.},
    type=dinge
}

% =====================================================================
%  Karteikarten  Vorlesung 3
% =====================================================================

\newglossaryentry{intronursprung}{
    name={Nenne die Modelle/Hypothesen zum Ursprung von Introns und erklaere sie kurz},
    description={introns-first (introns-early): Introns waren in fruehen Genen bereits vorhanden und gingen z.B. bei Prokaryoten sekundaer verloren. introns-late: fruehe Gene hatten keine Introns; diese wurden erst spaeter eingefuegt. Exon-Shuffling: neue Gene entstehen durch Rekombination/Neukombination von Exons, die haeufig fuer Protein(sub)domaenen kodieren.},
    type=dinge
}

\newglossaryentry{klassenrepetitiv}{
    name={Nenne die Hauptklassen repetitiver Sequenzen im eukaryotischen Genom},
    description={Satelliten-DNA / Minisatelliten (hochrepetitiv), Mikrosatelliten (kurze SSR), SINEs und LINEs (mittel-repetitiv, transponieren ueber RNA), Alu (ein SINE) sowie Genfamilien (duplizierte Gene mit ggf. differenzierter Funktion).},
    type=dinge
}


% =====================================================================
%  Karteikarten  Vorlesung 2
% =====================================================================
\newglossaryentry{baccontigs}{
    name={BAC-Contigs (Prinzip)},
    description={Das Genom wird in grosse DNA-Stuecke zerlegt, in BACs kloniert; anschliessend wird bestimmt, welche Stuecke sich ueberlappen. Aus den Ueberlappungen wird die urspruengliche Reihenfolge im Genom rekonstruiert (physikalische Karte).},
    type=dinge
}

\newglossaryentry{restriktionskartierung}{
    name={Restriktionskartierung (Prinzip)},
    description={DNA wird mit verschiedenen Restriktionsenzymen -- einzeln und paarweise -- verdaut, im Agarosegel getrennt und die Fragmentgroessen bestimmt. Aus dem Vergleich der Einzel- und Doppelverdaue werden die Fragmente zur Karte kombiniert.},
    type=dinge
}

\newglossaryentry{prokeuk}{
    name={Unterschiede Prokaryoten vs. Eukaryoten},
    description={Prokaryoten: kernlos, evolutionaer aelter, Gene in Operons (polycistronische mRNA), wenig/keine Introns, Transkription und Translation gekoppelt (Translation am nascenten Transkript). Eukaryoten: Zellkern, ausgepraegte Kompartimentierung (Endosymbiontentheorie), i.d.R. monocistronische mRNA, Gene mit Introns als Regel, mRNA-Prozessierung (Cap, Polyadenylierung), Transkription (Kern) und Translation (Cytoplasma) raeumlich getrennt, Mitose/Meiose.},
    type=dinge
}


\newglossaryentry{cistronisch}{
    name={Unterschiede mono- vs. polycistronische mRNA},
    description={Monocistronische mRNA (typisch eukaryotisch) kodiert fuer ein Protein. Polycistronische mRNA (typisch prokaryotisch, aus Operon) kodiert fuer mehrere Proteine.},
    type=dinge
}

\newglossaryentry{genzahlgenomgroesse}{
    name={Genzahl vs. Genomgroesse},
    description={Bei Bakterien und Archaeen sind Genzahl und Genomgroesse proportional. Bei Eukaryoten korreliert die Genzahl nicht gut mit der Genomgroesse oder der Komplexitaet des Organismus. Einzellige Eukaryoten liegen im Bereich grosser bakterieller Genome; Mehrzeller haben mehr Gene, aber ohne enge Korrelation zur Genomgroesse.},
    type=dinge
}

\newglossaryentry{organellengenom}{
    name={Eigenschaften von Organellengenomen},
    description={Meist ein einzelnes ringfoermiges DNA-Molekuel mit einzigartiger Sequenz (wenige Ausnahmen: lineare mtDNA bei einzelligen Eukaryoten). Mehrere Kopien pro Organelle und viele Organellen pro Zelle $\rightarrow$ viele Genomkopien pro Zelle; das Organellengenom kann als repetitive Sequenz betrachtet werden.},
    type=dinge
}

\newglossaryentry{genetvsphys}{
    name={Genetische vs. physikalische Karte},
    description={Genetische (Kopplungs-)Karte: Abstaende aus Rekombinationshaeufigkeiten (cM), abhaengig von beobachtbaren Markern. Physikalische Karte: Abstaende in Basenpaaren (z.B. aus Restriktionsschnittstellen, gemessen ueber Gelwanderung). Beide ordnen Loci, aber in unterschiedlichen Einheiten.},
    type=dinge
}

\newglossaryentry{annotationswege}{
    name={Von Sequenz zu Annotation (3 Wege)},
    description={(1) Ab initio (intrinsische Sequenzsignale). (2) Comparative Genomics (Konservierung zwischen Arten). (3) Datenbanksuchen/Alignments (Homologie zu bekannten Sequenzen).},
    type=dinge
}

\newglossaryentry{gesamtgenom}{
    name={Gesamt-Genom eines Eukaryoten (3 Kompartimente)},
    description={Besteht aus dem Nukleom (nDNA im Kern), dem Chondrom (mtDNA im Mitochondrium) und bei Pflanzen zusaetzlich dem Plastom (cpDNA im Chloroplasten).},
    type=dinge
}

\newglossaryentry{organellenimport}{
    name={Organellen kodieren nicht alle ihre Proteine, sondern ...},
    description={Organellengenome kodieren nur einen Teil der in der Organelle benoetigten Proteine. Die meisten urspruenglichen Organellengene wurden im Lauf der Evolution in den Kern verlagert; ihre Produkte werden im Cytoplasma synthetisiert und in die Organelle importiert.},
    type=dinge
}


